% SPDX-License-Identifier: AGPL-3.0-or-later
% Commercial license available
% © Concepts 1996-2026 Miroslav Sotek. All rights reserved.
% © Code 2020-2026 Miroslav Sotek. All rights reserved.
% ORCID: 0009-0009-3560-0851
% Contact: www.anulum.li | protoscience@anulum.li
% scpn-quantum-control -- Rust/VQE methodology manuscript draft

\documentclass[preprint,aps,physrev,superscriptaddress,floatfix,longbibliography]{revtex4-2}

\usepackage{amsmath,amssymb}
\usepackage{booktabs}
\usepackage{graphicx}
\usepackage{hyperref}
\usepackage{xcolor}
\usepackage{siunitx}
\usepackage{seqsplit}

\hypersetup{
    colorlinks=true,
    linkcolor=blue!60!black,
    citecolor=blue!60!black,
    urlcolor=blue!60!black,
}

\newcommand{\paperdoi}{Zenodo DOI: \href{https://doi.org/10.5281/zenodo.20382138}{10.5281/zenodo.20382138}}

\begin{document}
\emergencystretch=3em

\title{A Rust-accelerated, topology-informed workflow for Kuramoto--XY
quantum simulation and NISQ benchmarking}

\author{Miroslav \v{S}otek}
\email{protoscience@anulum.li}
\affiliation{ANULUM, Marbach SG, Switzerland}
\affiliation{ORCID: 0009-0009-3560-0851}

\date{\today\\\paperdoi}

\begin{abstract}
We describe the software and experimental workflow behind
\texttt{scpn-quantum-control}, an open-source framework for mapping
heterogeneous Kuramoto-type oscillator networks to XY spin Hamiltonians,
compiling topology-informed variational ansatze, benchmarking classical
and quantum simulation paths, and executing reproducible NISQ hardware
experiments. The framework combines Python/Qiskit orchestration with a
Rust/PyO3 acceleration layer for hot-path kernels such as coupling-matrix
construction, Hamiltonian assembly, trajectory features, hybrid
digital--analog partitioning, and protected-subspace diagnostics. We report
artefact-generated timings on three CPU sources (local workstation, ML350,
and Vertex AI \texttt{n1-standard-4}) and one Vertex AI Tesla T4 GPU run.
Fresh benchmark artefacts generated on 2026-05-05 show direct
Rust coupling-matrix speedups of $15.5$--$44.1\times$ at $n=16$
across the three self-contained CPU runs, exact checksum parity across
Python, Rust, and Go kernels, and a
three-seed $n=4$ VQE comparison in which the topology-informed ansatz
achieves a mean relative energy error of $0.455\,\%$, compared with
$6.136\,\%$ for \texttt{TwoLocal} and $7.522\,\%$ for
\texttt{EfficientSU2} at the same iteration budget. A follow-up
convergence artefact extends the same comparison to 320 COBYLA
evaluations and records seed-wise uncertainty: at $n=4$ and one
ansatz repetition the topology-informed family reaches a mean relative
error of $3.08\times 10^{-8}\,\%$, compared with $6.76\times10^{-4}\,\%$
for \texttt{TwoLocal} and $1.92\times10^{-3}\,\%$ for
\texttt{EfficientSU2}. A batched
state-vector expectation benchmark on Vertex T4 gives CUDA median
times of $0.176$--$0.905\,\mathrm{ms}$ for $n=8$--$11$, compared with
Vertex CPU medians of $3.557$--$12.094\,\mathrm{ms}$ for the same
workload. IBM ansatz-validation lanes were then executed at $n=4,6,8$ with
full computational-basis readout calibration. The original $n=4$ Marrakesh
lane preserves the topology-informed ordering, but the larger $n=6,8$
Fez/Marrakesh lanes rank \texttt{EfficientSU2} lowest in measured energy.
Additional $n=8$ seed-sensitivity lanes with seed 23 preserve that
larger-width ordering on both backends.
The hardware evidence is therefore a reproducibility and ansatz-sensitivity
check, not a backend-general topology-informed superiority claim. The paper is
methodological rather than a quantum-advantage claim: small-system exact
simulation remains classically tractable, and hardware results are used
to validate compilation, noise budgeting, and experiment reproducibility.
All code, benchmark commands, hardware raw counts, and derived artefacts are public at
\href{https://github.com/anulum/scpn-quantum-control}{\texttt{anulum/scpn-quantum-control}}.
\end{abstract}

\maketitle

\section{Introduction}

Quantum simulation papers often underreport the engineering layer that
connects a model Hamiltonian to executable circuits, raw-count archives,
analysis scripts, and reproducible performance claims. This omission is
especially problematic for NISQ studies, where small changes in ansatz
topology, transpilation, hardware calibration, readout processing, or
classical baseline implementation can dominate the advertised physics.

This manuscript presents the computational methodology used in the SCPN
Kuramoto--XY campaign~\cite{SotekSCPNSoftware}.
The contribution is
not a broad quantum-advantage result. Instead, it is a reproducible
pipeline for:
%
\begin{enumerate}
  \item compiling heterogeneous oscillator couplings into XY Hamiltonians,
  \item generating coupling-topology-informed ansatze and Trotter circuits,
  \item dispatching hot kernels to Rust while preserving Python/Qiskit
        usability,
  \item benchmarking classical, simulator, and hardware paths with
        command-level provenance, and
  \item archiving raw IBM hardware counts with integrity hashes and
        reproduction scripts.
\end{enumerate}

The immediate use case is Kuramoto-type synchronisation mapped to the
XY Hamiltonian
%
\begin{equation}
  H_{XY} =
  \sum_{i<j} K_{ij}\left(X_i X_j + Y_i Y_j\right)
  + \sum_i \omega_i Z_i,
\end{equation}
%
where $K_{ij}$ encodes the oscillator coupling topology and $\omega_i$
are heterogeneous natural frequencies. The same workflow also supports
VQE, Trotterised dynamics, readout/post-processing studies, DLA-parity
diagnostics, and hardware experiment packaging.

The paper is intentionally separated from the companion IBM hardware
preprint. The hardware preprint asks whether a parity-sector and
excitation-number correlated leakage signal is reproducible on real Heron
r2 hardware~\cite{SotekDLA2026,SotekDLAData2026}. The present manuscript asks a different question: whether
the software stack that generated, validated, and analysed those
experiments is itself reproducible, benchmarked, and reusable.

\section{Software architecture and artefact discipline}

The framework is organised as a Python front-end with a Rust/PyO3 engine
for hot paths. The Python layer provides model construction, Qiskit
circuits, hardware runners, analysis scripts, and documentation. The Rust
layer provides deterministic kernels for coupling construction, trajectory
features, Hamiltonian-adjacent operations, and selected diagnostics. This
division keeps the public API accessible while reducing overhead in loops
where Python object allocation would dominate.

The central reproducibility rule in this manuscript is that numerical
tables are not hand-authored. Benchmarks are generated by scripts under
\texttt{scripts/} and written to JSON/CSV artefacts under
\texttt{data/rust\_vqe\_methods/}. The manuscript cites those artefacts
directly. This prevents stale performance claims and makes later
reruns auditable when compilers, CPUs, cloud machines, or Python
dependencies change.

\begin{table}[tb]
\centering
\small
\begin{tabular}{lll}
\toprule
Layer & Representative modules & Role \\
\midrule
Bridge & \texttt{bridge/knm\_hamiltonian.py} & $K_{ij}\rightarrow H_{XY}$ \\
Ansatz & \texttt{control/structured\_ansatz.py} & topology-informed circuits \\
Phase & \texttt{phase/phase\_vqe.py} & VQE and dynamics workflows \\
Hardware & \texttt{hardware/runner.py} & simulator / IBM execution \\
Rust & \texttt{scpn\_quantum\_engine} & hot-path kernels \\
\bottomrule
\end{tabular}
\caption{High-level workflow components. The paper reports only
auditable functions with command-level provenance or committed test
coverage.}
\label{tab:architecture}
\end{table}

\section{Minimal working example}

The core user workflow is intentionally short: construct a Kuramoto--XY
coupling matrix, convert it into an operator, generate a topology-informed
ansatz, and then pass the same artefacts to a VQE or benchmark harness.
The following sketch is representative of the public API used by the
benchmark scripts:

\begin{verbatim}
from scpn_quantum_control.bridge.knm_hamiltonian import (
    OMEGA_N_16,
    build_knm_paper27,
    knm_to_ansatz,
    knm_to_hamiltonian,
)

n = 4
k_matrix = build_knm_paper27(n)
hamiltonian = knm_to_hamiltonian(k_matrix, OMEGA_N_16[:n])
ansatz = knm_to_ansatz(k_matrix, reps=2)
\end{verbatim}

The benchmark artefacts in \texttt{data/rust\_vqe\_methods/} are produced
by extending this workflow with fixed seeds, fixed optimiser budgets,
explicit machine metadata, and JSON/CSV output. This keeps the examples
close to the submitted hardware and methods data rather than presenting
toy APIs unrelated to the reported numbers.

\section{Rust kernel boundary}

The Rust layer is deliberately narrow. It accelerates deterministic kernels
that are called repeatedly from Python workflows, while Qiskit remains the
circuit, operator, and transpilation layer. For the coupling-matrix
benchmark, the accelerated operation has the conceptual signature

\begin{verbatim}
pub fn build_knm(n: usize, k0: f64, alpha: f64) -> Vec<f64>
\end{verbatim}

and returns a row-major dense representation of the exponentially decaying
coupling matrix. The benchmark harness compares this path against matched
Python and Go implementations using checksum parity before reporting timing
statistics. The reported speedups therefore refer to this kernel boundary,
not to the entire Qiskit workflow.

This distinction is an Amdahl's-law caveat rather than a weakness. For
example, reducing a $K_{ij}$ construction call from roughly
$10^{-2}\,\mathrm{ms}$ to $10^{-3}\,\mathrm{ms}$ is useful when the kernel is
called many times inside scans or packaging loops, but it does not remove the
dominant costs of an end-to-end VQE or hardware workflow. Qiskit circuit
construction and transpilation, optimiser evaluations, state-vector
expectation calculations, GPU/CPU dense linear algebra, queue latency, and QPU
execution can each dominate the wall time in different regimes. The software
claim is therefore bounded to reducing specific deterministic bottlenecks and
improving reproducibility, not to accelerating every workflow stage by the
reported kernel factor.

\section{Topology-informed ansatz design}

The topology-informed ansatz uses the support of $K_{ij}$ to decide which
qubit pairs receive entangling gates. This is a simple design principle:
do not spend variational capacity on edges absent from the physical
coupling graph. For sparse or structured oscillator networks, this can
reduce parameter count and circuit clutter relative to generic linear or
all-to-all ansatze.

The current benchmark harness regenerates ansatz construction and VQE
tables from source. For the representative $n=4$, two-repetition case,
the topology-informed ansatz uses fewer parameters than the generic
alternatives while retaining the same raw two-qubit-gate count as
\texttt{TwoLocal}. The generated VQE artefact then compares all three
ansatze over three random seeds with identical COBYLA iteration budgets.
The original comparison is deliberately reported as a fixed-budget score,
not as a converged variational optimum: the May 5 artefact records the
terminal energies after 80 COBYLA iterations per seed, but did not publish
a full energy-versus-iteration convergence trace. The 2026-05-20 extension closes
that specific documentation gap for the representative $n=4$ lane by
publishing checkpointed best-so-far energies, final parameters, elapsed-time
spread, and standard-error columns. It does not remove the broader
small-system limitation, but it distinguishes an 80-evaluation screening
result from a longer optimiser-trace result.

\begin{table}[tb]
\centering
\small
\begin{tabular}{lcccc}
\toprule
Ansatz & Params & Raw depth & 2q gates & Median VQE error \\
\midrule
$K_{ij}$-informed & 16 & 13 & 12 & $0.372\,\%$ \\
\texttt{TwoLocal} & 24 & 15 & 12 & $7.164\,\%$ \\
\texttt{EfficientSU2} & 24 & 11 & 6 & $7.708\,\%$ \\
\bottomrule
\end{tabular}
\caption{Generated $n=4$, two-repetition ansatz/VQE benchmark. Circuit
columns come from
\texttt{ansatz\_benchmark\_summary\_2026-05-05.json};
VQE errors come from
\texttt{vqe\_benchmark\_summary\_2026-05-05.json}
with three seeds and 80 COBYLA iterations per seed.}
\label{tab:ansatz}
\end{table}

A follow-up scaling artefact extends the circuit-size and truncation
diagnostics to $n=4,6,8,10,12$. Dense exact diagonalisation is used up to
$n=8$, and sparse \texttt{eigsh} ground-state references are used for
$n=10,12$. The generated sparse residual norms are
$1.12\times10^{-11}$ at $n=10$ and $8.70\times10^{-10}$ at $n=12$.
The corresponding MPS truncation diagnostic gives mid-chain entropies of
$0.9999$ bits at $n=10$ and $1.3143$ bits at $n=12$ for the canonical
Kuramoto--XY coupling matrix. These rows strengthen the scaling
diagnostic, but they remain ground-state reference data rather than a
large-system VQE optimisation theorem.

\begin{table}[tb]
\centering
\small
\begin{tabular}{lrrrr}
\toprule
Ansatz & Mean error & Std. err. & Best error & Mean time \\
\midrule
$K_{ij}$-informed & $3.08\times10^{-8}\,\%$ & $1.36\times10^{-8}\,\%$ & $1.07\times10^{-8}\,\%$ & $187.0\,\mathrm{ms}$ \\
\texttt{TwoLocal} & $6.76\times10^{-4}\,\%$ & $6.76\times10^{-4}\,\%$ & $5.59\times10^{-8}\,\%$ & $612.8\,\mathrm{ms}$ \\
\texttt{EfficientSU2} & $1.92\times10^{-3}\,\%$ & $1.69\times10^{-3}\,\%$ & $2.93\times10^{-5}\,\%$ & $618.7\,\mathrm{ms}$ \\
\bottomrule
\end{tabular}
\caption{Generated $n=4$, one-repetition VQE convergence extension from
\texttt{vqe\_convergence\_methods\_2026-05-20.json}. Each row uses three
seeds and 320 COBYLA evaluations. The table reports best-so-far relative
energy error against exact diagonalisation; the standard error is across
the three seeds.}
\label{tab:vqeconvergence}
\end{table}

\begin{table}[tb]
\centering
\small
\begin{tabular}{lrrrr}
\toprule
Ansatz & $n$ & Params & Transpiled depth & 2q gates \\
\midrule
$K_{ij}$-informed & 16 & 32 & 68 & 115 \\
\texttt{TwoLocal} & 16 & 64 & 72 & 120 \\
\texttt{EfficientSU2} & 16 & 64 & 25 & 15 \\
$K_{ij}$-informed & 20 & 40 & 84 & 163 \\
\texttt{TwoLocal} & 20 & 80 & 88 & 190 \\
\texttt{EfficientSU2} & 20 & 80 & 29 & 19 \\
\bottomrule
\end{tabular}
\caption{Generated construction/transpilation scaling extension from
\texttt{ansatz\_scaling\_error\_bars\_2026-05-20.csv}. This table is
not a dense VQE energy claim at $n=16$--20; it is the memory-safe
large-$n$ resource lane requested by the methods review.}
\label{tab:largenscaling}
\end{table}

\section{Generated CPU benchmark artefacts}

Table~\ref{tab:performance} lists conservative local performance facts generated
by the new benchmark harnesses. These are opportunistic timings on a shared
workstation, not final publication-grade timing claims; CPU load from other
processes was not pinned or isolated. The older internal claim of a
$5401\times$ Hamiltonian-construction speedup is not used here because it
has not yet been reproduced by the present harness with a matched baseline.
The fresh coupling-matrix benchmark shows useful speedups at $n=4$--$32$,
but not at $n=64$; the manuscript therefore reports bounded kernel-level
acceleration rather than a universal speedup.

To make the comparison extensible, we added a multi-language
coupling-matrix harness. It currently reports Python/NumPy, direct
Rust/PyO3, Julia, and Go timings where local runtimes are available.
Runtimes without a stable matched $K_{ij}$ kernel are intentionally omitted
from the result table rather than listed as failed benchmarks.

\begin{table*}[tb]
\centering
\small
\begin{tabular}{lrrr}
\toprule
Task & $n$ & Median time & Note \\
\midrule
\texttt{build\_knm}, Python & 4 & $0.00627\,\mathrm{ms}$ & reference \\
\texttt{build\_knm}, Rust & 4 & $0.00106\,\mathrm{ms}$ & $5.89\times$ \\
\texttt{build\_knm}, Python & 16 & $0.01028\,\mathrm{ms}$ & reference \\
\texttt{build\_knm}, Rust & 16 & $0.00277\,\mathrm{ms}$ & $3.71\times$ \\
\texttt{build\_knm}, Python & 32 & $0.01148\,\mathrm{ms}$ & reference \\
\texttt{build\_knm}, Rust & 32 & $0.00828\,\mathrm{ms}$ & $1.39\times$ \\
Dense $H_{XY}$ & 8 & $0.11705\,\mathrm{ms}$ & Hermitian error 0 \\
\emph{Timing context} & -- & -- & shared workstation; non-isolated \\
\bottomrule
\end{tabular}
\caption{Generated opportunistic shared-workstation core benchmark table from
\texttt{rust\_core\_benchmark\_summary\_2026-05-05.json}. Python/Rust
coupling matrices match to the configured tolerance for all tested
sizes $n\in\{4,8,16,32,64\}$. The dense $H_{XY}$ row reports the maximum
absolute Hermiticity residual, which is zero for the generated matrix.}
\label{tab:performance}
\end{table*}

\begin{table}[tb]
\centering
\small
\begin{tabular}{lcc}
\toprule
Language path & $n=16$ median & Status \\
\midrule
Python/NumPy & $0.01339\,\mathrm{ms}$ & ok \\
Rust/PyO3 & $0.00330\,\mathrm{ms}$ & ok \\
Julia & $0.00509\,\mathrm{ms}$ & ok \\
Go & $0.00402\,\mathrm{ms}$ & ok \\
\bottomrule
\end{tabular}
\caption{Generated multi-language $K_{ij}$ construction comparison at
$n=16$ from
\texttt{multilang\_knm\_benchmark\_summary\_2026-05-05.json}.
The same shared-workstation caveat applies; this table is a comparator
framework, not yet a final language shootout.}
\label{tab:multilang}
\end{table}

The self-contained cross-machine harness was then executed on three
independent CPU sources: the local workstation \texttt{aaarthuus}, the
ML350 server \texttt{god-of-the-math}, and a Vertex AI
\texttt{n1-standard-4} worker. This harness deliberately avoids importing
the project package so that a secondary machine can run the same
Python/Rust/Go kernels before the full project environment is installed.
All three runs report matching checksums for the three language kernels,
which makes the timings meaningful as implementation comparisons rather
than unrelated algorithms.

\paragraph*{Cross-machine $K_{ij}$ construction comparison.}
The generated 16-oscillator cross-machine benchmark is kept inline rather than
as a floating table so the machine rows remain adjacent to the explanatory
text.
\begin{center}
\small
\begin{tabular}{lrrr}
\toprule
Machine & Python median & Rust median & Rust/Python \\
\midrule
Local & $0.06631\,\mathrm{ms}$ & $0.00150\,\mathrm{ms}$ & $44.1\times$ \\
ML350 & $0.07620\,\mathrm{ms}$ & $0.00493\,\mathrm{ms}$ & $15.5\times$ \\
Vertex CPU & $0.04705\,\mathrm{ms}$ & $0.00260\,\mathrm{ms}$ & $18.1\times$ \\
\bottomrule
\end{tabular}
\end{center}
The rows are generated from the remote $K_{ij}$ benchmark JSONs and the
combined methods benchmark summary JSON dated 2026-05-05. Each row uses the
self-contained benchmark harness with 1000 iterations and matching
Python/Rust/Go checksums; the Go timings are retained in the JSON artefacts but
omitted from this compact manuscript table. The reported machines had non-zero
load averages during measurement, so the table supports robustness of the
ordering, not isolated microbenchmark claims.

\section{GPU-assisted validation benchmark}

GPU acceleration is useful for a different part of the workflow than the
Rust kernels. The $K_{ij}$ construction kernel is small and allocation
sensitive; moving it to a GPU would mostly benchmark transfer overhead.
The GPU-relevant workload is batched classical validation: evaluating
many state-vector energies or VQE candidate states against the same
dense Hamiltonian. We therefore added a separate GPU harness for batched
expectation values,
\begin{equation}
  E_b = \langle \psi_b | H_{XY} | \psi_b \rangle,
  \qquad b=1,\ldots,B,
\end{equation}
with $B=64$ random normalised complex state vectors.

Local GPU timings are not reported as valid because the local NVIDIA
stack was inconsistent during this session: \texttt{nvidia-smi} failed
with an NVML driver/library mismatch, and CUDA initialisation in the
local CUDA-enabled PyTorch environment failed with error 804. Instead,
the valid GPU artefact is the Vertex AI Tesla T4 run using
\texttt{pytorch/pytorch:2.5.1-cuda12.4-cudnn9-runtime}.

\begin{table}[tb]
\centering
\small
\begin{tabular}{rrrr}
\toprule
$n$ & Hilbert dim. & Vertex CPU median & T4 CUDA median \\
\midrule
8 & 256 & $11.121\,\mathrm{ms}$ & $0.176\,\mathrm{ms}$ \\
10 & 1024 & $3.557\,\mathrm{ms}$ & $0.290\,\mathrm{ms}$ \\
11 & 2048 & $12.094\,\mathrm{ms}$ & $0.905\,\mathrm{ms}$ \\
\bottomrule
\end{tabular}
\caption{Generated batched expectation benchmark from
\texttt{gpu\_benchmark\_summary\_vertex\_t4\_2026-05-05.json}. The
workload uses batch size 64 and 20 timing repeats. The CPU column is the
NumPy path inside the same Vertex worker; the CUDA column is the PyTorch
CUDA path on one Tesla T4 using the
\texttt{pytorch/pytorch:2.5.1-cuda12.4-cudnn9-runtime} container.}
\label{tab:gpu}
\end{table}

These results support a hybrid implementation strategy. Rust is used for
low-latency CPU kernels and deterministic orchestration hot paths; GPU
acceleration is used for batched dense linear algebra in classical
validation and VQE scoring. QPU time is reserved for hardware-noise
questions that cannot be answered by classical simulation alone.

\section{Hardware workflow and reproducibility}

The same repository contains the raw-count hardware workflow used for the
DLA-parity paper. That campaign is not the central result of the present
methodology paper, but it demonstrates why the software architecture is
useful: each promoted hardware dataset carries raw counts, circuit
metadata, IBM job identifiers, SHA256 hashes, and reproducer scripts. This
paper should cite the DLA-parity preprint as an application case study,
while keeping the main methodological claim focused on the toolchain and
benchmarking discipline.

For the methods ansatz claim itself, the 2026-05-20 revision adds a
bounded hardware validation lane on \texttt{ibm\_marrakesh}. The lane used
the same $n=4$ ansatz families and fixed seed as the convergence artefact,
measured the $Z$, $X$, and $Y$ global bases for each ansatz, and included a
full 16-state computational-basis readout calibration. The job
\texttt{d86u48op0eas73dlfqs0} executed 25 circuits at 1024 shots each on
physical qubits $[7,17,6,8]$. The calibration assignment matrix was
well-conditioned, with condition number $1.059$.

\begin{table}[tb]
\centering
\small
\begin{tabular}{lrr}
\toprule
Ansatz & Raw energy & Full-readout mitigated energy \\
\midrule
$K_{ij}$-informed & $-6.1938$ & $-6.2021$ \\
\texttt{TwoLocal} & $-6.0733$ & $-6.0805$ \\
\texttt{EfficientSU2} & $-6.0461$ & $-6.0539$ \\
\bottomrule
\end{tabular}
\caption{IBM Marrakesh ansatz-validation lane from the 2026-05-20 analysis
JSON listed in the data-availability section. The result is a small hardware
ordering check, not a hardware convergence proof. It supports the narrower
statement that the topology-informed ansatz energy ordering survived one
pinned IBM backend/layout execution and full 16-state readout-mitigation
replay.}
\label{tab:hardwareansatz}
\end{table}

To address the small-width limitation directly, we also submitted the same
ansatz-validation machinery at $n=6$ and $n=8$ on \texttt{ibm\_fez} and
\texttt{ibm\_marrakesh}. Each larger lane measured the same three ansatz
families and included the complete $2^n$ computational-basis readout
calibration. These runs validate the packaging, retrieval, and readout
inversion path at larger width, but they do not preserve the $n=4$
topology-informed ranking. Table~\ref{tab:hardwareansatzlarge} reports the
full-readout-mitigated energies. In all four larger lanes,
\texttt{EfficientSU2} has the lowest measured energy. This is a useful negative
methods result: the topology-informed ansatz remains a good small-system design
signal, but the larger hardware lanes show that ansatz ranking is sensitive to
width, backend, layout, and noise.

\begin{table}[tb]
\centering
\small
\begin{tabular}{lllrrr}
\toprule
Backend & $n$ & Seed & $K_{ij}$-informed & \texttt{TwoLocal} & \texttt{EfficientSU2} \\
\midrule
\texttt{ibm\_fez} & 6 & 11 & $-9.1343$ & $-7.9625$ & $-10.0840$ \\
\texttt{ibm\_fez} & 8 & 11 & $-6.6274$ & $-3.6932$ & $-9.7473$ \\
\texttt{ibm\_fez} & 8 & 23 & $-6.7392$ & $-5.1422$ & $-10.0636$ \\
\texttt{ibm\_marrakesh} & 6 & 11 & $-10.0442$ & $-9.9373$ & $-10.3118$ \\
\texttt{ibm\_marrakesh} & 8 & 11 & $-11.3559$ & $-10.2083$ & $-11.9303$ \\
\texttt{ibm\_marrakesh} & 8 & 23 & $-11.4752$ & $-10.6883$ & $-11.9551$ \\
\bottomrule
\end{tabular}
\caption{Larger-width IBM ansatz-validation lanes submitted on 2026-05-20.
Entries are full-readout-mitigated energy estimates from the retrieved raw
counts. The seed-23 repeat at $n=8$ preserves the larger-width
\texttt{EfficientSU2}-lowest ordering on both backends. The result is
deliberately reported as ansatz-sensitivity evidence: the larger hardware
lanes validate the workflow at $n=6,8$ but do not promote the
topology-informed ansatz as hardware-general superior.}
\label{tab:hardwareansatzlarge}
\end{table}

\begin{table}[tb]
\centering
\small
\begin{tabular}{lp{0.62\linewidth}}
\toprule
Workflow component & Interpretation \\
\midrule
$K_{ij}$ construction & Rust/PyO3 kernel acceleration reduces repeated
coupling-graph construction cost during sweeps and artefact regeneration. \\
Hamiltonian/ansatz export & Python/Qiskit object construction remains part of
the cost and limits end-to-end speedup. \\
Optimiser evaluations & COBYLA and state-vector evaluations dominate VQE
screening once the Hamiltonian is built. \\
Transpilation/provider execution & Hardware lanes are governed by circuit
depth, backend routing, shot budget, queueing, and calibration windows. \\
\bottomrule
\end{tabular}
\caption{End-to-end placement of the Rust kernel in the methods workflow. The
kernel speedups are real and useful, but they are deliberately reported as
workflow-component gains rather than full-stack VQE speedups.}
\label{tab:endtoendbreakdown}
\end{table}

The larger-width hardware lanes also give a useful negative design rule. The
topology-informed ansatz is the best family in the original small $n=4$
screening, but \texttt{EfficientSU2} gives the lowest measured energy in the
$n=6,8$ hardware lanes and in the seed-23 $n=8$ repeats. A plausible
mechanism is that the topology-informed circuit encodes non-native graph
structure whose routed two-qubit gates and calibration-window sensitivity
outweigh its variational inductive bias on current heavy-hex hardware. The
paper therefore promotes the ansatz path as a reproducible compiler and
benchmarking method, not as a universal hardware ansatz-selection rule.

\section{Limitations}

The benchmark results are deliberately bounded. First, the CPU timings
are opportunistic shared-machine timings. They record machine metadata
and load averages, but do not pin CPU affinity, clock governor, cache
state, or thermal state. They are sufficient to document the pipeline and
compare implementation ordering; they are not a final language shootout.

Second, the VQE comparison is small. The $n=4$ topology-informed ansatz
result is a useful methods signal, but it should not be interpreted as a
universal optimisation theorem. The original fixed 80-iteration COBYLA budget
is a controlled screening budget, not proof of convergence; the 320-evaluation
extension improves traceability for the representative case but still covers
only a small Hamiltonian and three seeds. Larger ansatz matrices, more seeds,
and additional Hamiltonian families remain necessary before promoting a broad
ansatz-selection rule. The IBM observable comparisons in
Tables~\ref{tab:hardwareansatz} and~\ref{tab:hardwareansatzlarge} are hardware
sanity checks for the methods workflow, not backend-general ansatz-superiority
evidence.

Third, the GPU benchmark is a classical-validation workload, not a claim
that the full Qiskit circuit-construction path is GPU accelerated. Circuit
construction and transpilation remain mostly Python/object-graph
workloads. The GPU contribution is batched dense linear algebra.

Finally, the hardware experiments that motivated this pipeline remain
small-N NISQ experiments. The software stack improves reproducibility and
throughput, but it does not by itself establish quantum advantage.
At the time of the reported hardware campaigns, authenticated live QPU access
for the programme was limited to IBM Quantum hardware. The package includes
provider-neutral execution and artefact-packaging surfaces for additional
gate-model, trapped-ion, neutral-atom, photonic, annealing, analogue, and
simulator targets, but these routes should be read as readiness infrastructure,
not as live non-IBM hardware evidence. The intended next use of the methods
stack is to rerun the same preregistered payloads on additional hardware
families when compute allocations and provider access become available.

\section{Data and code availability}

The public repository is
\href{https://github.com/anulum/scpn-quantum-control}{\texttt{https://github.com/anulum/scpn-quantum-control}}.
The methods-paper benchmark artefacts are under
\path{data/rust_vqe_methods/}. The primary generated summaries used
in this manuscript are:
%
\begin{itemize}
  \item \path{rust_core_benchmark_summary_2026-05-05.json},
  \item \path{ansatz_benchmark_summary_2026-05-05.json},
  \item \path{vqe_benchmark_summary_2026-05-05.json},
  \item \path{multilang_knm_benchmark_summary_2026-05-05.json},
  \item \path{remote_knm_benchmark_aaarthuus_local_2026-05-05.json},
  \item \path{remote_knm_benchmark_ml350_2026-05-05.json},
  \item \path{remote_knm_benchmark_vertex_n1_standard_4_2026-05-05.json},
  \item \path{gpu_benchmark_summary_vertex_t4_2026-05-05.json},
  \item \path{ansatz_scaling_tn_summary_2026-05-05.json},
  \item \path{ansatz_scaling_summary_2026-05-05.csv},
  \item \path{tn_truncation_summary_2026-05-05.csv},
  \item \path{combined_methods_benchmark_summary_2026-05-05.json},
  \item \path{vqe_convergence_methods_2026-05-20.json},
  \item \path{vqe_convergence_aggregate_2026-05-20.csv},
  \item \path{ansatz_scaling_error_bars_2026-05-20.csv},
  \item \path{ansatz_ibm_validation_readiness_ibm_marrakesh_20260520T163234Z.json},
  \item \path{ansatz_ibm_validation_submission_ibm_marrakesh_20260520T163234Z.json},
  \item \path{ansatz_ibm_validation_raw_counts_ibm_marrakesh_20260520T163812Z.json},
  \item \path{ansatz_ibm_validation_analysis_ibm_marrakesh_20260520T163812Z.json},
  \item \path{ansatz_ibm_validation_analysis_ibm_fez_20260520T195028Z.json},
  \item \path{ansatz_ibm_validation_analysis_ibm_fez_20260520T195040Z.json},
  \item \path{ansatz_ibm_validation_analysis_ibm_marrakesh_20260520T195051Z.json},
  \item \path{ansatz_ibm_validation_analysis_ibm_marrakesh_20260520T195102Z.json},
  \item \path{ansatz_ibm_validation_analysis_ibm_fez_20260520T201721Z.json}, and
  \item \path{ansatz_ibm_validation_analysis_ibm_marrakesh_20260520T201741Z.json}.
\end{itemize}
%
The latest combined summary hashes at the time of this draft are
\texttt{\seqsplit{593330a1dd19f495b899be1031ebe3dd4caa07171053aa376c2f761e557c1428}}
for JSON and
\texttt{\seqsplit{e69b94df590ff06708b3b21245864f74c3df630b514254526dc6c4af3fe24c2f}}
for CSV. The corresponding generator scripts are
\path{scripts/benchmark_rust_core_methods.py},
\path{scripts/benchmark_ansatz_methods.py},
\path{scripts/benchmark_vqe_methods.py},
\path{scripts/benchmark_multilang_knm_methods.py},
\path{scripts/benchmark_remote_knm_machine.py},
\path{scripts/benchmark_gpu_methods.py}, and
\path{scripts/summarise_rust_vqe_method_artifacts.py}. The 2026-05-20
extension scripts are
\path{scripts/benchmark_vqe_convergence_methods.py},
\path{scripts/submit_methods_ansatz_ibm_validation.py}, and
\path{scripts/retrieve_methods_ansatz_ibm_validation.py}.

\section{Conclusion}

We have presented a conservative methodology paper around
\texttt{scpn-quantum-control}: a Rust-accelerated, topology-informed
workflow for Kuramoto--XY quantum simulation, VQE benchmarking, and NISQ
hardware reproducibility. The strongest near-term contribution is not an
unqualified speedup headline, but a transparent engineering stack that
connects physical coupling graphs to circuits, benchmark artefacts,
hardware raw counts, and reproducible analysis. The generated artefacts
show that Rust gives consistent low-latency CPU benefits for the
coupling-matrix kernel across three CPU sources, that topology-informed
ansatz design improves the tested small VQE case, and that Vertex T4 GPU
acceleration is useful for batched classical expectation evaluation.

\begin{thebibliography}{9}

\bibitem{SotekSCPNSoftware}
M.~\v{S}otek,
\emph{\texttt{scpn-quantum-control}: Quantum-Native SCPN Phase Dynamics
and Control}, version 0.9.7,
Zenodo DOI:~\href{https://doi.org/10.5281/zenodo.18821929}{10.5281/zenodo.18821929} (2026).

\bibitem{SotekDLA2026}
M.~\v{S}otek,
``Replicated hardware observation of parity-sector and excitation-number
correlated decoherence asymmetry in the XY Hamiltonian,''
Zenodo preprint DOI:~\href{https://doi.org/10.5281/zenodo.20382000}{10.5281/zenodo.20382000} (2026).

\bibitem{SotekDLAData2026}
M.~\v{S}otek,
\emph{DLA parity hardware validation package for heterogeneous
Kuramoto--XY circuits},
Zenodo dataset and public repository artefacts (2026),
Zenodo record 19098792, \url{https://zenodo.org/records/19098792}.

\bibitem{Qiskit}
Qiskit contributors,
\emph{Qiskit: An Open-source Framework for Quantum Computing}.

\bibitem{Peruzzo2014}
A.~Peruzzo \emph{et al.},
``A variational eigenvalue solver on a photonic quantum processor,''
Nature Communications \textbf{5}, 4213 (2014).

\bibitem{Kandala2017}
A.~Kandala \emph{et al.},
``Hardware-efficient variational quantum eigensolver for small molecules
and quantum magnets,''
Nature \textbf{549}, 242--246 (2017).

\bibitem{Kuramoto1984}
Y.~Kuramoto,
\emph{Chemical Oscillations, Waves, and Turbulence}
(Springer, Berlin, 1984).

\end{thebibliography}

\end{document}
